\documentclass[12pt]{article}
\usepackage[a4paper, margin=1in]{geometry}
\usepackage{amsmath, amssymb, booktabs, setspace}
\setstretch{1.2}
\parskip=0.7em



\title{Tutorial  --The Heckscher--Ohlin Model}
\author{ECON 308 -- International Economic Relations}
\date{Fall 2025}

\begin{document}
\maketitle



\noindent
Two countries, \textbf{Home (Canada)} and \textbf{Foreign (Mexico)}, produce two goods — \textbf{Computers (C)} and \textbf{Textiles (T)} — using two factors of production: \textbf{Capital (K)} and \textbf{Labor (L)}.  
Both countries share identical technologies and preferences. The production functions for each good are:

\[
Q_C = (K_C)^{0.6}(L_C)^{0.4}, \quad Q_T = (K_T)^{0.3}(L_T)^{0.7}
\]

\noindent
The total factor endowments are:

\begin{center}
\begin{tabular}{lcc}
\toprule
\textbf{Country} & \textbf{Capital (K)} & \textbf{Labor (L)} \\
\midrule
Canada (Home) & 600 & 300 \\
Mexico (Foreign) & 300 & 600 \\
\bottomrule
\end{tabular}
\end{center}

\noindent
Computers are \textbf{capital-intensive}, while Textiles are \textbf{labor-intensive}.  
Both goods are produced under perfect competition and constant returns to scale.

%============================================================
\section*{1. Factor Abundance}

\noindent\textbf{Question:}
\begin{enumerate}
    \item Compute the \( K/L \) ratio for each country. Which country is capital-abundant and which is labor-abundant?
    \item Suppose both countries have the same technology. If Canada doubles both its capital and labor, does its factor abundance classification change?
\end{enumerate}

\noindent\textbf{Solution and Explanation:}

\[
(K/L)_{Canada} = \frac{600}{300} = 2.0, \quad (K/L)_{Mexico} = \frac{300}{600} = 0.5
\]

\noindent
\textbf{Definition:}  
A country is \emph{capital-abundant} if it has a higher ratio of capital to labor than another country; it is \emph{labor-abundant} if the opposite holds.

\noindent
Thus, Canada (2.0) is capital-abundant, and Mexico (0.5) is labor-abundant.

\noindent
If both capital and labor in Canada double, its new ratio is \(1200/600 = 2.0\); therefore, its classification does not change.  
Factor abundance depends on \textbf{relative proportions}, not absolute quantities.

%============================================================
\section*{2. Factor Intensity}

\noindent\textbf{Question:}
\begin{enumerate}
    \item Using:
    \[
    K_C = 300, \; L_C = 100; \quad K_T = 150, \; L_T = 350
    \]
    compute the \(K/L\) ratio in each industry.
    \item Derive the cost-minimizing capital-labor ratio for a Cobb–Douglas production function and explain how the wage–rental ratio (\(w/r\)) affects input choice.
\end{enumerate}

\noindent\textbf{Solution and Explanation:}

\[
(K/L)_C = \frac{300}{100} = 3.0, \quad (K/L)_T = \frac{150}{350} = 0.43
\]

\noindent
Thus, Computers use more capital per worker — they are \textbf{capital-intensive}.

\noindent
\textbf{Definition:}  
An industry is \emph{capital-intensive} if its \(K/L\) ratio exceeds that of another industry at given factor prices.

\noindent
For a Cobb–Douglas function \(Q = K^\alpha L^{1-\alpha}\), cost minimization yields:
\[
\frac{K}{L} = \frac{\alpha}{1-\alpha}\frac{w}{r}
\]
This shows that when the wage–rental ratio \(w/r\) increases (labor becomes expensive), firms substitute toward capital, increasing \(K/L\).
\newpage
%============================================================
\section*{3. Autarky Equilibrium and Relative Prices}

\noindent\textbf{Question:}
\begin{enumerate}
    \item Suppose:
    \[
    (w/r)_{Canada} = 1.5, \quad (w/r)_{Mexico} = 3.0
    \]
    Compute the implied \(K/L\) ratios for Computers and Textiles in each country using the cost-minimization condition.
    \item Which good is relatively cheaper in each country under autarky? Explain.
    \item Qualitatively describe the relative supply curves for both countries and indicate the direction of comparative advantage.
\end{enumerate}

\noindent\textbf{Solution and Explanation:}

\textbf{Step 1: Compute Capital-Labor Ratios}

For Computers (\(\alpha = 0.6\)):
\[
(K/L)_C = \frac{0.6}{0.4} \times \frac{w}{r} = 1.5 \times \frac{w}{r}
\]

For Textiles (\(\alpha = 0.3\)):
\[
(K/L)_T = \frac{0.3}{0.7} \times \frac{w}{r} = 0.43 \times \frac{w}{r}
\]

Substituting each country's \(w/r\):

\begin{center}
\begin{tabular}{lcc}
\toprule
 & \((K/L)_C\) & \((K/L)_T\) \\
\midrule
Canada & $1.5 \times 1.5 = 2.25$ & $0.43 \times 1.5 = 0.65$ \\
Mexico & $1.5 \times 3.0 = 4.5$ & $0.43 \times 3.0 = 1.29$ \\
\bottomrule
\end{tabular}
\end{center}

\noindent
In both countries, Computers remain more capital-intensive.
\newpage
\textbf{Step 2: Relative Prices}

- In Canada, capital is cheap (\(w/r\) low) → Computers (capital-intensive) are relatively cheaper.  
- In Mexico, labor is cheap (\(w/r\) high) → Textiles (labor-intensive) are relatively cheaper.  

\[
(P_C/P_T)_{Canada} < (P_C/P_T)_{Mexico}
\]

\textbf{Step 3: Comparative Advantage}

Canada has a comparative advantage in Computers; Mexico in Textiles.  
Hence, the relative supply curve of Computers for Canada lies farther to the right.

%============================================================
\section*{4. Trade Pattern and Terms of Trade}

\noindent\textbf{Question:}
\begin{enumerate}
    \item Predict each country’s exports and imports once trade opens.
    \item If autarky prices are:
    \[
    (P_C/P_T)_{Canada} = 1.2, \quad (P_C/P_T)_{Mexico} = 2.2
    \]
    and world price after trade is \( (P_C/P_T)^W = 1.8 \), determine how each country’s \textbf{terms of trade} change and the welfare implications.
\end{enumerate}

\noindent\textbf{Solution and Explanation:}

\textbf{Trade Pattern (Heckscher–Ohlin Theorem):}  
Each country exports the good that uses its abundant factor intensively.
\begin{itemize}
    \item Canada (capital-abundant) exports Computers and imports Textiles.
    \item Mexico (labor-abundant) exports Textiles and imports Computers.
\end{itemize}

\textbf{Terms of Trade:}  
\[
\text{ToT} = \frac{P_{export}}{P_{import}}
\]
- For Canada, \(P_C/P_T\) rises from 1.2 to 1.8 → terms of trade improve.  
- For Mexico, \(P_C/P_T\) falls from 2.2 to 1.8 → terms of trade deteriorate.  

\textbf{Welfare Effects:}  
Even though factor owners are affected differently, both countries gain overall since they can now consume beyond their production possibilities.



\end{document}
